\documentclass[12pt,a4paper]{report}

\usepackage{graphicx}
\usepackage{amsmath}
\usepackage{geometry}
\usepackage{setspace}
\usepackage{enumitem}
\usepackage{hyperref}
\usepackage{booktabs}
\usepackage{array}
\usepackage{pgfgantt}
\usepackage{listings}
\usepackage{xcolor}
\usepackage{float}

\hypersetup{
    colorlinks=true,
    linkcolor=black,
    citecolor=blue,
    urlcolor=cyan
}

\geometry{margin=1in}
\onehalfspacing

\usepackage{titlesec}

% Remove "Chapter" word and use numeric headings
\titleformat{\chapter}[hang]
{\normalfont\LARGE\bfseries}
{\thechapter.}{1em}{}

\titlespacing*{\chapter}{0pt}{-10pt}{20pt}

% ---- Table column types ----
\newcolumntype{L}[1]{>{\raggedright\arraybackslash}p{#1}}
\newcolumntype{C}[1]{>{\centering\arraybackslash}p{#1}}

% ---- Verilog listing style ----
\lstdefinestyle{verilog}{
    language=Verilog,
    basicstyle=\ttfamily\footnotesize,
    keywordstyle=\color{blue}\bfseries,
    commentstyle=\color{green!50!black}\itshape,
    stringstyle=\color{red!70!black},
    numbers=left,
    numberstyle=\tiny\color{gray},
    numbersep=5pt,
    frame=single,
    breaklines=true,
    tabsize=4,
    showstringspaces=false,
    captionpos=b
}

% ---- Console output style ----
\lstdefinestyle{console}{
    basicstyle=\ttfamily\footnotesize,
    frame=single,
    breaklines=true,
    backgroundcolor=\color{gray!10},
    captionpos=b
}


\begin{document}
\pagenumbering{gobble}

%================ TITLE PAGE ==================
\begin{titlepage}
\centering

% INSERT TU LOGO PNG HERE (NAME FILE: Logo/TU_Logo.png)
% \includegraphics[width=0.25\textwidth]{Logo/TU_Logo.png}\\[0.5cm]

{\Large TRIBHUVAN UNIVERSITY}\\
{\Large INSTITUTE OF ENGINEERING}\\
{\Large PULCHOWK CAMPUS}\\[1cm]

{\bfseries\LARGE A MID-TERM REPORT}\\[0.5cm]
{\Large ON}\\[0.5cm]
{\bfseries\LARGE DESIGN AND IMPLEMENTATION OF \\[0.3cm] RISC-V PIPELINING}\\[2cm]

SUBMITTED BY:\\[0.3cm]
\textbf{Aayush Gelal (079BEI002)}\\
\textbf{Adesh Pandey (079BEI004)}\\
\textbf{Bishesh Paudel (079BEI016)}\\[1cm]

SUBMITTED TO:\\
Department of Electronics \& Computer Engineering\\[1cm]

March 2026
\end{titlepage}

\clearpage
\pagenumbering{roman}

%================ ACKNOWLEDGEMENT ==================
\chapter*{Acknowledgement}
\addcontentsline{toc}{chapter}{Acknowledgement}

We would like to express our sincere gratitude to the Department of Electronics and
Computer Engineering for providing the academic environment and laboratory resources
necessary to carry out this project. We are deeply grateful to our teachers and the
departmental faculty for their continued guidance and technical oversight throughout the
design and implementation phases. Their expertise in RISC-V microarchitecture and
hardware design principles has been invaluable in shaping the direction of this work.

We also acknowledge the support of our peers and seniors, whose collaborative
discussions have helped refine our implementation approach. Furthermore, we appreciate
the resources provided by the global RISC-V open-source community, which have been
essential in understanding the industry-standard specifications required for this work.
As we approach the hardware synthesis and FPGA deployment phase, we look forward to
the continued mentorship of our faculty and the department.

%================ ABSTRACT ==================
\chapter*{Abstract}
\addcontentsline{toc}{chapter}{Abstract}

This mid-term report documents the progress made toward the design and implementation
of a five-stage pipelined processor based on the RV32I RISC-V instruction set
architecture. At the current stage of the project, the processor design has progressed
through two major phases: a fully functional single-cycle implementation supporting
the complete RV32I instruction set, followed by a successful conversion to a five-stage
pipelined architecture with comprehensive hazard handling.

The pipelined processor implements all base RV32I instructions including arithmetic,
logical, shift, comparison, load, store, branch, jump, and upper immediate operations.
The pipeline incorporates a data forwarding unit that resolves read-after-write hazards
from both the EX/MEM and MEM/WB stages, a load-use hazard detection unit that inserts
pipeline stalls when forwarding alone is insufficient, and a control hazard mechanism
that flushes incorrectly fetched instructions on taken branches and jumps.

All pipeline hazard mechanisms have been verified through simulation using Icarus
Verilog, with a comprehensive test program that exercises EX/MEM forwarding, MEM/WB
forwarding, load-use stalling, store forwarding, R-type forwarding chains, and branch
flushing with forwarded operands. All test cases pass with correct results. The
remaining work focuses on synthesizing the design on an FPGA platform for hardware
validation.

\tableofcontents
\addcontentsline{toc}{chapter}{List of Figures}
\listoffigures
\addcontentsline{toc}{chapter}{List of Tables}
\listoftables

%================ LIST OF ABBREVIATIONS ==================
\chapter*{List of Abbreviations}
\addcontentsline{toc}{chapter}{List of Abbreviations}

\begin{tabular}{ll}
ALU   & Arithmetic Logic Unit \\
CPI   & Cycles Per Instruction \\
CPU   & Central Processing Unit \\
EX    & Execute \\
FPGA  & Field Programmable Gate Array \\
HDL   & Hardware Description Language \\
ID    & Instruction Decode \\
IF    & Instruction Fetch \\
ISA   & Instruction Set Architecture \\
MEM   & Memory Access \\
NOP   & No Operation \\
PC    & Program Counter \\
RAW   & Read After Write \\
RISC  & Reduced Instruction Set Computer \\
RTL   & Register Transfer Level \\
RV32I & 32-bit RISC-V Base Integer ISA \\
VCD   & Value Change Dump \\
WB    & Write Back \\
\end{tabular}

\clearpage
\pagenumbering{arabic}
\setcounter{page}{1}

%================ CHAPTER 1: INTRODUCTION ==================
\chapter{Introduction}

\section{Background}
Processor architecture has evolved significantly since the early days of computing,
transitioning from simple, hardwired control units to highly parallel and deeply
pipelined designs. At the core of this evolution lies the instruction set architecture,
which defines the interface between software and hardware. Traditional commercial ISAs
are often proprietary, limiting transparency and hindering academic exploration.

RISC-V is an open, modular instruction set architecture developed to address these
limitations. Its clean design, extensibility, and royalty-free nature make it an ideal
candidate for educational processor design projects. Implementing a RISC-V processor
provides students with practical exposure to fundamental architectural concepts such as
instruction decoding, datapath control, and memory interfacing.

\section{Objectives}
The primary objectives of this project are as follows:

\begin{itemize}
    \item To design and implement a 32-bit processor compliant with the RISC-V RV32I
          instruction set architecture.
    \item To implement a classical five-stage pipelined datapath with hazard handling.
    \item To support all RV32I instruction formats and operations.
    \item To verify functional correctness through simulation.
    \item To synthesize and implement the processor on an FPGA platform.
\end{itemize}

\section{Scope}
The scope of this project is limited to the design, simulation, and FPGA implementation
of a 32-bit pipelined RISC-V processor core. The processor supports all RV32I integer
instructions, including arithmetic, logical, shift, comparison, load, store, branch,
and jump instructions.

Advanced features such as floating-point extensions, vector processing, cache memory,
virtual memory, operating system support, and complex peripheral interfaces are beyond
the scope of this work. The focus is strictly on the processor core, pipeline
operation, and correct instruction execution.

\section{Work Completed to Date}
As of this mid-term report, the following major milestones have been achieved:

\begin{itemize}
    \item Design and implementation of all individual hardware modules in Verilog HDL:
          ALU, Control Unit, Register File, Program Counter, Immediate Generator,
          Instruction Memory, and Data Memory.
    \item Full integration and verification of a single-cycle RV32I CPU supporting
          all base integer instructions.
    \item Successful conversion of the single-cycle design to a five-stage pipelined
          architecture (\texttt{top.v}) with four pipeline registers (IF/ID, ID/EX,
          EX/MEM, MEM/WB).
    \item Implementation of a data forwarding unit supporting EX/MEM and MEM/WB
          forwarding paths.
    \item Implementation of a load-use hazard detection unit with pipeline stalling.
    \item Implementation of control hazard handling via pipeline flushing for branches
          and jumps.
    \item Comprehensive pipeline hazard verification through a 46-instruction test
          program with all tests passing.
\end{itemize}

The remaining work consists of synthesizing the pipelined design on an FPGA platform
for hardware validation and performance evaluation.


%================ CHAPTER 2: MODULE DESIGN ==================
\chapter{Module-Level Design and Implementation}

This chapter describes each Verilog module designed and implemented during the
project. All modules are written using synthesizable RTL constructs compatible with
Icarus Verilog simulation and FPGA synthesis tools. Table~\ref{tab:modules}
summarizes the project source files.

\begin{table}[h!]
\centering
\caption{Project Source File Summary}
\label{tab:modules}
\begin{tabular}{L{3.2cm} L{2.4cm} L{6.8cm}}
    \toprule
    \textbf{Module} & \textbf{File} & \textbf{Function} \\
    \midrule
    ALU              & \texttt{alu.v}       & Arithmetic and logical operations \\
    Control Unit     & \texttt{control.v}   & Opcode decoding and control signal generation \\
    Register File    & \texttt{regfile.v}   & 32$\times$32-bit register storage \\
    Program Counter  & \texttt{pc.v}        & Instruction address register with enable \\
    PC Adder         & \texttt{pc\_adder.v} & PC + 4 increment logic \\
    Immediate Gen.   & \texttt{imm\_gen.v}  & Sign-extended immediate extraction \\
    Instruction Mem. & \texttt{imem.v}      & Read-only instruction storage (64 words) \\
    Data Memory      & \texttt{dmem.v}      & Read/write data storage (64 words) \\
    Top-Level CPU    & \texttt{top.v}       & 5-stage pipeline with hazard handling \\
    \bottomrule
\end{tabular}
\end{table}

\section{Arithmetic Logic Unit (ALU)}
The ALU module (\texttt{alu.v}) implements all arithmetic and logical operations
required by the RV32I instruction set. It accepts two 32-bit operands and a 4-bit
control selector, producing a 32-bit result and a zero flag. The ALU is implemented
as a purely combinational module using a \texttt{case} statement. The supported
operations are listed in Table~\ref{tab:aluops}.

\begin{table}[h!]
\centering
\caption{ALU Operations and Control Codes}
\label{tab:aluops}
\begin{tabular}{C{2.4cm} L{3.6cm} L{5.4cm}}
    \toprule
    \textbf{Control} & \textbf{Operation} & \textbf{Description} \\
    \midrule
    \texttt{4'b0000} & ADD  & 32-bit addition \\
    \texttt{4'b0001} & SUB  & 32-bit subtraction \\
    \texttt{4'b0010} & AND  & Bitwise AND \\
    \texttt{4'b0011} & OR   & Bitwise OR \\
    \texttt{4'b0100} & XOR  & Bitwise XOR \\
    \texttt{4'b0101} & SLT  & Signed set-less-than \\
    \texttt{4'b0110} & SLL  & Shift left logical \\
    \texttt{4'b0111} & SRL  & Shift right logical \\
    \texttt{4'b1000} & SLTU & Unsigned set-less-than \\
    \texttt{4'b1001} & SRA  & Shift right arithmetic \\
    \bottomrule
\end{tabular}
\end{table}

\section{Control Unit}
The Control Unit (\texttt{control.v}) decodes the 7-bit opcode, 3-bit \texttt{funct3},
and 7-bit \texttt{funct7} fields to generate all control signals required for
instruction execution. The control signals include \texttt{reg\_write},
\texttt{alu\_src}, \texttt{alu\_control}, \texttt{mem\_write}, \texttt{result\_src},
\texttt{branch}, \texttt{jump}, \texttt{jalr}, and \texttt{alu\_src1}.

The \texttt{result\_src} signal is a 2-bit selector: \texttt{2'b00} selects the ALU
result, \texttt{2'b01} selects memory read data, and \texttt{2'b10} selects PC+4 for
link instructions (JAL, JALR). The \texttt{alu\_src1} signal selects the ALU first
operand source: \texttt{2'b00} for register data, \texttt{2'b01} for the PC value
(AUIPC), and \texttt{2'b10} for zero (LUI). Table~\ref{tab:control} summarizes the
control signal values for each instruction class.

\begin{table}[h!]
\centering
\caption{Control Unit Signal Truth Table}
\label{tab:control}
\begin{tabular}{L{2.0cm} C{1.5cm} C{1.2cm} C{1.2cm} C{1.2cm} C{1.5cm} C{1.2cm} C{1.2cm}}
    \toprule
    \textbf{Type} & \textbf{Opcode} & \textbf{RegWr} & \textbf{ALU Src} &
    \textbf{MemWr} & \textbf{ResSrc} & \textbf{Branch} & \textbf{Jump} \\
    \midrule
    R-Type    & 0110011 & 1 & 0 & 0 & 00 & 0 & 0 \\
    I-Type    & 0010011 & 1 & 1 & 0 & 00 & 0 & 0 \\
    Load      & 0000011 & 1 & 1 & 0 & 01 & 0 & 0 \\
    Store     & 0100011 & 0 & 1 & 1 & -- & 0 & 0 \\
    Branch    & 1100011 & 0 & 0 & 0 & -- & 1 & 0 \\
    JAL       & 1101111 & 1 & 0 & 0 & 10 & 0 & 1 \\
    JALR      & 1100111 & 1 & 1 & 0 & 10 & 0 & 0 \\
    LUI       & 0110111 & 1 & 1 & 0 & 00 & 0 & 0 \\
    AUIPC     & 0010111 & 1 & 1 & 0 & 00 & 0 & 0 \\
    \bottomrule
\end{tabular}
\end{table}

\section{Register File}
The Register File (\texttt{regfile.v}) implements 32 general-purpose 32-bit registers.
It provides two asynchronous read ports and one synchronous write port. Register
\texttt{x0} is hardwired to zero: reads from address zero always return
\texttt{32'b0}, regardless of any prior write.

\section{Program Counter and PC Adder}
The Program Counter module (\texttt{pc.v}) maintains the current instruction address.
It features a synchronous reset and an enable signal used by the hazard detection unit
to stall the PC during load-use hazards. The PC Adder (\texttt{pc\_adder.v}) is a
simple combinational module that adds 4 to the current PC for sequential fetch.

\section{Immediate Generator}
The Immediate Generator (\texttt{imm\_gen.v}) extracts and sign-extends the immediate
field from the instruction based on the opcode. It supports all six RV32I immediate
formats:

\begin{itemize}
    \item \textbf{I-Type} (ADDI, loads, JALR): \texttt{\{\{20\{instr[31]\}\}, instr[31:20]\}}
    \item \textbf{S-Type} (stores): \texttt{\{\{20\{instr[31]\}\}, instr[31:25], instr[11:7]\}}
    \item \textbf{B-Type} (branches): 13-bit with implicit LSB zero
    \item \textbf{J-Type} (JAL): 21-bit with implicit LSB zero
    \item \textbf{U-Type} (LUI, AUIPC): \texttt{\{instr[31:12], 12'b0\}}
\end{itemize}

\section{Instruction and Data Memory}
The Instruction Memory (\texttt{imem.v}) is a 64-word read-only memory initialized
from a hexadecimal file (\texttt{program.hex}) using \texttt{\$readmemh}. The Data
Memory (\texttt{dmem.v}) is a 64-word synchronous-write, asynchronous-read memory.
Both memories use word-aligned addressing via \texttt{a[31:2]}.


%================ CHAPTER 3: SINGLE-CYCLE TO PIPELINE ==================
\chapter{From Single-Cycle to Pipelined Architecture}

This chapter describes the evolution of the processor from a single-cycle design to a
five-stage pipelined architecture, which constitutes the core technical contribution
of this project phase.

\section{Single-Cycle Design}
The initial implementation connected all modules into a single-cycle datapath where
each instruction completes in one clock cycle. The single-cycle design was developed
incrementally:

\begin{enumerate}
    \item Basic arithmetic instructions (ADD, ADDI) were implemented first.
    \item Data memory operations (LW, SW) were added with the Data Memory module.
    \item Branch instructions (BEQ, BNE, BLT, BGE, BLTU, BGEU) were implemented
          with combinational branch condition evaluation.
    \item Jump instructions (JAL, JALR) were added with PC+4 link writeback.
    \item The complete RV32I instruction set was finalized, including LUI, AUIPC,
          and all remaining I-type, R-type, and shift operations.
\end{enumerate}

While functionally correct, the single-cycle design has a fundamental performance
limitation: the clock period must accommodate the longest instruction path (load
instruction through fetch, decode, ALU, memory, and writeback), leaving faster
instructions idle for the majority of the cycle.

\section{Five-Stage Pipeline Architecture}
To improve instruction throughput, the single-cycle datapath was partitioned into five
pipeline stages with registers inserted between each pair of adjacent stages:

\begin{enumerate}
    \item \textbf{Instruction Fetch (IF)}: The PC drives the instruction memory. The
          fetched instruction and PC value are written to the IF/ID register.
    \item \textbf{Instruction Decode (ID)}: The control unit decodes the instruction.
          The register file is read and the immediate is generated. All values and
          control signals are written to the ID/EX register.
    \item \textbf{Execute (EX)}: The ALU performs the operation. Branch conditions are
          evaluated. Forwarded operands are selected. Results are written to the
          EX/MEM register.
    \item \textbf{Memory Access (MEM)}: Data memory is read or written. The result
          and control signals propagate to the MEM/WB register.
    \item \textbf{Write Back (WB)}: The result multiplexer selects between ALU
          output, memory data, or PC+4, and writes back to the register file.
\end{enumerate}

\section{Pipeline Register Design}
Four pipeline registers are implemented as synchronous registers clocked on the
positive edge. Each register stores both datapath values and control signals needed
by subsequent stages. Table~\ref{tab:pipereg} lists the contents of each pipeline
register.

\begin{table}[H]
\centering
\caption{Pipeline Register Contents}
\label{tab:pipereg}
\begin{tabular}{L{2.0cm} L{10.2cm}}
    \toprule
    \textbf{Register} & \textbf{Contents} \\
    \midrule
    IF/ID  & Instruction (32b), PC (32b), PC+4 (32b) \\
    ID/EX  & PC, PC+4, rs1 data, rs2 data, immediate, rd address, rs1/rs2
             addresses, funct3, and all control signals (reg\_write, alu\_src,
             mem\_write, branch, jump, jalr, result\_src, alu\_src1, alu\_control) \\
    EX/MEM & ALU result, store data (forwarded rs2), PC+4, rd address,
             reg\_write, mem\_write, result\_src \\
    MEM/WB & ALU result, memory read data, PC+4, rd address, reg\_write,
             result\_src \\
    \bottomrule
\end{tabular}
\end{table}

% ---- FIGURE: Five-Stage Pipeline Architecture ----
% Uncomment and provide your own diagram:
% \begin{figure}[H]
% \centering
% \includegraphics[width=0.85\textwidth]{pictures/5-stage.png}
% \caption{Five-Stage RISC-V Pipeline Architecture (IF, ID, EX, MEM, WB)}
% \label{fig:pipeline}
% \end{figure}


%================ CHAPTER 4: HAZARD HANDLING ==================
\chapter{Pipeline Hazard Detection and Resolution}

Pipelined execution introduces hazards that can lead to incorrect results if not
properly managed. This chapter describes the three hazard handling mechanisms
implemented in the processor.

\section{Data Forwarding Unit}
The forwarding unit detects RAW (read-after-write) data dependencies between
instructions in different pipeline stages. Two forwarding control signals,
\texttt{forward\_a} and \texttt{forward\_b}, independently control the source of
each ALU operand:

\begin{itemize}
    \item \texttt{2'b00}: No forwarding; use the value from the ID/EX register.
    \item \texttt{2'b10}: Forward from EX/MEM stage (1-instruction gap).
    \item \texttt{2'b01}: Forward from MEM/WB stage (2-instruction gap).
\end{itemize}

The EX/MEM path has priority over the MEM/WB path, since a more recent instruction's
result takes precedence when both match. The forwarding logic checks that the
destination register is not \texttt{x0} and that the producing instruction has
\texttt{reg\_write} asserted.

For JAL/JALR instructions in the EX/MEM stage, the forwarded value is PC+4 rather
than the ALU result, since these instructions write the return address to the
destination register.

\section{Load-Use Hazard Detection}
When a load instruction is in the EX stage and the immediately following instruction
in the ID stage reads the same register, forwarding alone cannot resolve the
dependency because the load data is not yet available. The hazard detection unit
asserts a \texttt{stall} signal when:

\begin{enumerate}
    \item The instruction in the ID/EX register has \texttt{result\_src == 2'b01}
          (indicating a load).
    \item The destination register is not \texttt{x0}.
    \item The destination register matches either the \texttt{rs1} or \texttt{rs2}
          field of the instruction in the IF/ID register.
\end{enumerate}

When \texttt{stall} is asserted:
\begin{itemize}
    \item The PC is frozen (enable deasserted).
    \item The IF/ID register holds its current value.
    \item A bubble (NOP) is inserted into the ID/EX register by clearing all control
          signals.
\end{itemize}

After one stall cycle, the load data becomes available and is forwarded normally from
the MEM/WB stage.

\section{Control Hazard Handling}
Branch and jump instructions are resolved in the EX stage. When a branch is taken or
a jump instruction is detected, the \texttt{flush} signal is asserted:

\[
\texttt{flush} = \texttt{ex\_branch\_taken} \;\|\; \texttt{id\_ex\_jump} \;\|\; \texttt{id\_ex\_jalr}
\]

The flush signal clears the IF/ID register (inserting a NOP instruction
\texttt{0x00000013}) and clears the ID/EX register (inserting a bubble). The PC is
simultaneously updated to the branch target (PC + immediate) or JALR target
(ALU result \& \texttt{0xFFFFFFFE}).

Branch conditions are evaluated in the EX stage using the ALU result and \texttt{funct3}:

\begin{table}[H]
\centering
\caption{Branch Condition Evaluation}
\label{tab:branch}
\begin{tabular}{C{2.0cm} L{3.0cm} L{5.5cm}}
    \toprule
    \textbf{funct3} & \textbf{Instruction} & \textbf{Condition (taken when)} \\
    \midrule
    \texttt{3'b000} & BEQ  & ALU zero flag is set \\
    \texttt{3'b001} & BNE  & ALU zero flag is clear \\
    \texttt{3'b100} & BLT  & ALU result[0] is set (SLT) \\
    \texttt{3'b101} & BGE  & ALU result[0] is clear \\
    \texttt{3'b110} & BLTU & ALU result[0] is set (SLTU) \\
    \texttt{3'b111} & BGEU & ALU result[0] is clear \\
    \bottomrule
\end{tabular}
\end{table}

\section{Hazard Interaction Summary}
Table~\ref{tab:hazard_summary} summarizes how the pipeline responds to each hazard
type.

\begin{table}[H]
\centering
\caption{Hazard Resolution Summary}
\label{tab:hazard_summary}
\begin{tabular}{L{3.5cm} L{2.8cm} L{5.0cm}}
    \toprule
    \textbf{Hazard Type} & \textbf{Mechanism} & \textbf{Pipeline Action} \\
    \midrule
    RAW (1-cycle gap)    & EX/MEM forwarding  & Forward result to ALU input \\
    RAW (2-cycle gap)    & MEM/WB forwarding  & Forward result to ALU input \\
    Load-use             & Stall + forward    & 1-cycle stall, then MEM/WB forward \\
    Taken branch / jump  & Flush              & Clear IF/ID and ID/EX, redirect PC \\
    \bottomrule
\end{tabular}
\end{table}


%================ CHAPTER 5: SIMULATION ==================
\chapter{Simulation and Verification}

\section{Software Tools Used}

\subsection{Icarus Verilog}
Icarus Verilog is used as the primary compiler and simulator for Verilog HDL. It
supports synthesizable RTL constructs and provides accurate event-driven simulation.
All modules and the complete pipelined processor are compiled and simulated using
this tool.

\subsection{GTKWave}
GTKWave is employed as a waveform visualization and debugging tool. The simulation
generates a Value Change Dump (\texttt{.vcd}) file that is analyzed in GTKWave to
observe pipeline register transitions, forwarding path activations, stall cycles,
and flush events across clock cycles.

\section{Module-Level Testbenches}
Individual testbenches were written and executed for the ALU and Register File modules
to verify correct functionality in isolation before integration.

\begin{table}[H]
\centering
\caption{Module-Level Testbench Summary}
\label{tab:testbenches}
\begin{tabular}{L{2.4cm} L{2.8cm} L{5.8cm} C{1.6cm}}
    \toprule
    \textbf{Module} & \textbf{Testbench} & \textbf{Tests Performed} & \textbf{Result} \\
    \midrule
    ALU           & \texttt{alu\_tb.v}     & ADD, SUB, AND operations    & Pass \\
    Register File & \texttt{regfile\_tb.v} & Write/read-back, x0 guard  & Pass \\
    \bottomrule
\end{tabular}
\end{table}

\section{Pipeline Hazard Test Program}
The primary verification of the pipelined processor is performed through a
comprehensive test program (\texttt{program.hex}) consisting of 46 instructions
organized into six targeted test cases. Each test case exercises a specific pipeline
hazard scenario.

\subsection{Test 1: EX/MEM Forwarding (Back-to-Back Dependency)}
Three consecutive instructions create immediate data dependencies:
\begin{lstlisting}[style=console]
addi x1, x0, 5       # x1 = 5
addi x2, x1, 1       # x2 = x1 + 1 = 6  (forward x1 from EX/MEM)
addi x3, x2, 1       # x3 = x2 + 1 = 7  (forward x2 from EX/MEM)
\end{lstlisting}

\subsection{Test 2: MEM/WB Forwarding (2-Instruction Gap)}
A NOP separates the producing and consuming instructions, testing the MEM/WB
forwarding path:
\begin{lstlisting}[style=console]
addi x4, x0, 10      # x4 = 10
nop                   # 1-instruction gap
addi x5, x4, 4       # x5 = x4 + 4 = 14  (forward x4 from MEM/WB)
\end{lstlisting}

\subsection{Test 3: Load-Use Hazard (Stall + Forward)}
A load immediately followed by a dependent instruction triggers the stall mechanism:
\begin{lstlisting}[style=console]
sw   x1, 0(x0)       # mem[0] = 5
...                   # (NOPs to let store complete)
lw   x6, 0(x0)       # x6 = mem[0] = 5
addi x7, x6, 1       # x7 = x6 + 1 = 6   (stall 1 cycle, then forward)
\end{lstlisting}

\subsection{Test 4: R-Type Forwarding Chain}
Back-to-back R-type instructions with register-to-register forwarding:
\begin{lstlisting}[style=console]
add  x1, x1, x2      # x1 = 5 + 6 = 11   (forward both operands)
add  x2, x1, x1      # x2 = 11 + 11 = 22 (forward x1 from EX/MEM)
\end{lstlisting}

\subsection{Test 5: Store with Forwarded Data}
A store instruction uses a recently computed value, verifying forwarding into the
store data path:
\begin{lstlisting}[style=console]
addi x8, x0, 20      # x8 = 20
sw   x8, 0(x0)       # mem[0] = 20  (store forwarded x8)
...
lw   x9, 0(x0)       # x9 = mem[0] = 20  (verify store)
\end{lstlisting}

\subsection{Test 6: Branch with Forwarding}
A taken branch uses forwarded operands for condition evaluation:
\begin{lstlisting}[style=console]
addi x10, x0, 5      # x10 = 5
addi x11, x0, 5      # x11 = 5
beq  x10, x11, +12   # taken (forward x10 and x11, flush 2 instrs)
addi x12, x0, 1      # SKIPPED (flushed)
addi x12, x0, 1      # SKIPPED (flushed)
addi x12, x0, 0      # x12 = 0 (branch target)
\end{lstlisting}

\section{Simulation Results}
The pipelined processor testbench (\texttt{top\_tb.v}) executes the test program for
600 time units (60 clock cycles) and then reads the final register values. The
complete simulation output is shown below:

\begin{lstlisting}[style=console,caption={Complete Simulation Output},label={lst:simout}]
=== Pipeline Hazard Tests ===
--- EX/MEM Forwarding (back-to-back) ---
x1  = 11 (expect 11, final after test4)
x2  = 22 (expect 22, final after test4)
x3  = 7  (expect 7,  chain fwd)
--- MEM/WB Forwarding (2-gap) ---
x5  = 14 (expect 14, x4+4)
--- Load-Use Hazard (stall + forward) ---
x6  = 5  (expect 5,  loaded)
x7  = 6  (expect 6,  load-use fwd)
--- Store Forwarding ---
x8  = 20 (expect 20)
x9  = 20 (expect 20, store-then-load)
--- Branch with Forwarding ---
x10 = 5  (expect 5)
x11 = 5  (expect 5)
x12 = 0  (expect 0,  branch taken)
\end{lstlisting}

All twelve register values match their expected results, confirming correct operation
of:
\begin{itemize}
    \item EX/MEM data forwarding (back-to-back and chained)
    \item MEM/WB data forwarding (2-instruction gap)
    \item Load-use hazard detection and stalling
    \item Store data forwarding
    \item Branch condition evaluation with forwarded operands
    \item Pipeline flushing on taken branches
\end{itemize}

\begin{table}[H]
\centering
\caption{Pipeline Verification Results Summary}
\label{tab:results}
\begin{tabular}{C{1.0cm} L{4.8cm} C{2.0cm} C{2.0cm} C{1.4cm}}
    \toprule
    \textbf{Test} & \textbf{Hazard Tested} & \textbf{Expected} & \textbf{Observed} & \textbf{Result} \\
    \midrule
    1 & EX/MEM forwarding      & x3 = 7   & x3 = 7   & Pass \\
    2 & MEM/WB forwarding      & x5 = 14  & x5 = 14  & Pass \\
    3 & Load-use stall         & x7 = 6   & x7 = 6   & Pass \\
    4 & R-type fwd chain       & x2 = 22  & x2 = 22  & Pass \\
    5 & Store forwarding       & x9 = 20  & x9 = 20  & Pass \\
    6 & Branch + flush         & x12 = 0  & x12 = 0  & Pass \\
    \bottomrule
\end{tabular}
\end{table}


%================ CHAPTER 6: SUPPORTED INSTRUCTIONS ==================
\chapter{Supported Instruction Set}

The processor supports all base RV32I instructions, ensuring full compliance with the
32-bit RISC-V specification. Table~\ref{tab:instructions} lists all supported
instructions organized by format type.

\begin{table}[H]
\centering
\caption{Complete RV32I Instruction Set Supported by the Processor}
\label{tab:instructions}
\begin{tabular}{L{2.2cm} L{2.0cm} L{7.2cm}}
    \toprule
    \textbf{Format} & \textbf{Opcode} & \textbf{Instructions} \\
    \midrule
    R-Type  & 0110011 & ADD, SUB, AND, OR, XOR, SLL, SRL, SRA, SLT, SLTU \\
    I-Type  & 0010011 & ADDI, ANDI, ORI, XORI, SLLI, SRLI, SRAI, SLTI, SLTIU \\
    Load    & 0000011 & LW \\
    Store   & 0100011 & SW \\
    Branch  & 1100011 & BEQ, BNE, BLT, BGE, BLTU, BGEU \\
    JAL     & 1101111 & JAL \\
    JALR    & 1100111 & JALR \\
    LUI     & 0110111 & LUI \\
    AUIPC   & 0010111 & AUIPC \\
    \bottomrule
\end{tabular}
\end{table}

The six RV32I instruction formats (R, I, S, B, U, J) are fully handled by the
Immediate Generator and Control Unit. All instruction encodings conform to the RISC-V
specification as defined in the unprivileged ISA manual \cite{riscv}.


%================ CHAPTER 7: REMAINING WORK ==================
\chapter{Remaining Work: FPGA Implementation}

With the pipelined processor fully implemented and verified through simulation, the
remaining work focuses on FPGA synthesis and hardware validation.

\section{FPGA Synthesis Strategy}
The processor is implemented using synthesizable Verilog HDL and targeted for
deployment on an FPGA platform. The design follows a modular approach with separate
Verilog modules for each functional unit, facilitating synthesis tool optimization.
Key synthesis tasks include:

\begin{itemize}
    \item Mapping the Verilog design onto FPGA configurable logic blocks, flip-flops,
          and block RAM resources.
    \item Applying clock constraints and performing static timing analysis to
          determine the maximum operating frequency.
    \item Evaluating resource utilization (LUTs, registers, memory blocks) and
          comparing with similar RISC-V implementations.
\end{itemize}

\section{Hardware Validation}
FPGA-based validation will complement RTL simulation by confirming correct processor
operation under real hardware conditions. The validation plan includes:

\begin{itemize}
    \item Programming the FPGA with the synthesized bitstream.
    \item Loading test programs into instruction memory.
    \item Observing processor behavior using on-board LEDs, switches, or a UART
          debugging interface.
    \item Verifying consistency between simulation results and hardware behavior.
\end{itemize}

\section{Implementation Considerations}
Several considerations guide the FPGA implementation phase:

\begin{itemize}
    \item \textbf{Synthesizability}: All Verilog constructs used in the design are
          compatible with FPGA synthesis tools. Non-synthesizable constructs
          (\texttt{\$readmemh}, \texttt{\$display}) are confined to simulation
          testbenches.
    \item \textbf{Timing Closure}: Pipeline stages are balanced to minimize critical
          path delays and achieve reliable operation.
    \item \textbf{Memory Initialization}: Instruction memory will be initialized
          using FPGA vendor-specific mechanisms (e.g., block RAM initialization
          files) to replace the simulation-only \texttt{\$readmemh}.
\end{itemize}


%================ CHAPTER 8: TIMELINE ==================
\chapter{Project Timeline}

Figure~\ref{fig:gantt_chart} shows the updated project Gantt chart. Phases 1 and 2
(single-cycle design and pipelined architecture) have been completed. The project now
enters Phase 3 (FPGA synthesis and hardware testing).

\begin{figure}[H]
\centering

\begin{ganttchart}[
    expand chart=\textwidth,
    vgrid,
    hgrid,
    y unit title=0.6cm,
    y unit chart=0.6cm,
    time slot format=simple,
    title/.style={draw=none, fill=none},
    title label font=\bfseries\small,
    bar height=0.6,
    bar label font=\scriptsize,
    group label font=\bfseries\scriptsize,
    bar label node/.append style={
        text width=3cm,
        align=center
    }
]{1}{60}

% ---- Time Labels ----
\gantttitle{January}{15}
\gantttitle{February}{28}
\gantttitle{March}{17} \\

% ---- Phase 1: COMPLETED ----
\ganttgroup{Phase 1: Single-Cycle (Done)}{1}{15} \\
\ganttbar[bar/.style={fill=green!55,draw=none}]{Module Design}{1}{10} \\
\ganttbar[bar/.style={fill=green!55,draw=none}]{Simulation \& Verify}{8}{15} \\

% ---- Phase 2: COMPLETED ----
\ganttgroup{Phase 2: Pipeline (Done)}{16}{30} \\
\ganttbar[bar/.style={fill=green!55,draw=none}]{Pipeline Registers}{16}{22} \\
\ganttbar[bar/.style={fill=green!55,draw=none}]{Hazard Unit}{21}{27} \\
\ganttbar[bar/.style={fill=green!55,draw=none}]{Pipeline Verification}{27}{30} \\

% ---- Phase 3: NEXT ----
\ganttgroup{Phase 3: FPGA Implementation}{31}{45} \\
\ganttbar[bar/.style={fill=orange!60,draw=none}]{Synthesis \& Timing}{31}{37} \\
\ganttbar[bar/.style={fill=orange!60,draw=none}]{Board Testing}{37}{45} \\

% ---- Phase 4: PLANNED ----
\ganttgroup{Phase 4: Reporting}{46}{60} \\
\ganttbar[bar/.style={fill=orange!60,draw=none}]{Final Documentation}{46}{55} \\

\end{ganttchart}

\caption{Updated Gantt Chart (Green = Completed, Orange = Remaining)}
\label{fig:gantt_chart}
\end{figure}


\clearpage

%================ REFERENCES ==================
\renewcommand{\bibname}{References}

\begin{thebibliography}{99}

\bibitem{patterson}
D. A. Patterson and J. L. Hennessy,
\textit{Computer Organization and Design: The Hardware/Software Interface},
5th ed., Morgan Kaufmann, 2014.

\bibitem{harris}
D. M. Harris and S. L. Harris,
\textit{Digital Design and Computer Architecture: RISC-V Edition},
Morgan Kaufmann, 2021.

\bibitem{asanovic}
K. Asanovi\'c and D. A. Patterson,
``Instruction Sets Should Be Free: The Case for RISC-V,''
EECS Department, University of California, Berkeley, Technical Report, 2014.

\bibitem{hennessy}
J. L. Hennessy and D. A. Patterson,
``A New Golden Age for Computer Architecture,''
\textit{Communications of the ACM},
vol. 62, no. 2, pp. 48--60, 2019.

\bibitem{riscv}
A. Waterman, K. Asanovi\'c, et al.,
\textit{The RISC-V Instruction Set Manual, Volume I: Unprivileged ISA},
RISC-V International, 2019.

\bibitem{pong}
P. Pong and R. C. C. Cheung,
``FPGA Implementation of a Pipelined RISC Processor,''
\textit{IEEE Transactions on Very Large Scale Integration (VLSI) Systems},
vol. 22, no. 3, pp. 605--618, 2014.

\end{thebibliography}

\end{document}
